\documentclass[12pt]{article}
\usepackage[utf8]{inputenc}
\usepackage[T1]{fontenc}
\usepackage{hyperref}
\title{Lang'e Language Resources for K-12 Curriculum Development}
\author{Compiled by Claude Sonnet 5.5 (claude-sonnet-5-5) for Andr\'as Kornai}
\date{}
\begin{document}
\maketitle

\section{Introduction}
Lang'e (ISO 639-3 yne; Glottocode lang1333) is a Loloish (Tibeto-Burman) language of Yunnan, China. It is spoken in about 12 villages of the Lang'e cluster in Changhai Township, Yongsheng County. Estimates are around 2,500 speakers (Wikipedia, undated) and about 2,200 (Ethnologue, via Joshua Project). The search found very little educational material. Joshua Project notes that little has been written about the group.

\section{Primary Educational Resources}
No primers, textbooks or school programmes were found.

\section{Linguistic and Reference Materials}
\begin{itemize}
\item Wikipedia article on Lang'e, with classification and location: \url{https://en.wikipedia.org/wiki/Lang%27e_language}
\end{itemize}
No dictionary or grammar was found in the search results.

\section{Religious and Community Resources}
\begin{itemize}
\item Joshua Project profile of the Lang'e people: \url{https://joshuaproject.net/people_groups/20604}
\item Joshua Project language page (yne): \url{https://joshuaproject.net/languages/yne}
\item Global Recordings Network language page (lists the alternate name La'u; the search reported no audio recordings): \url{https://globalrecordings.net/en/language/yne}
\end{itemize}

\section{Digital and Audio Resources}
No audio, video or software was found. The Global Recordings Network page lists no recordings for this language.

\section{Curriculum Development Context}
The language is reported as in everyday use within the community and not known to be taught in schools. A curriculum would need new documentation work with the community, starting with orthography and basic word lists. Nothing found indicates existing teaching materials.
\end{document}